\documentclass[12pt]{article}
\usepackage[utf8]{inputenc}
\usepackage[T1]{fontenc}
\usepackage{hyperref}
\title{Tapirap\'e Language Resources for K-12 Curriculum Development}
\author{Compiled by Claude Sonnet 5 (claude-sonnet-5) for Andr\'as Kornai}
\date{}
\begin{document}
\maketitle

\section{Introduction}
Tapirap\'e (Apy\~awa; ISO 639-3: taf; Glottocode: tapi1254) is a Tupi-Guaran\'i language spoken near the Araguaia River in Mato Grosso and Tocantins, Brazil. Wikipedia gives about 950 speakers (2020), versus about 560 in the reference description. It is taught in community schools. Search turned up academic studies rather than ready-made classroom materials.

\section{Primary Educational Resources}
No published Tapirap\'e textbook or primer was found in search results.
\begin{itemize}
\item Study of the bilingual school in the community (Rev. Bras. Ling. Aplicada): \url{https://www.scielo.br/j/rbla/a/xLNpDdSsQ6fW5F3B7fpwwHp/?lang=en}
\item Study of the Tapi'it\~awa State Indigenous School and Tapirap\'e neologisms: \url{https://periodicos.unb.br/index.php/lcapital/article/download/8598/7157/14864}
\item Article on modern foreign language at the Tapi'it\~awa school: \url{https://periodicos.cfs.ifmt.edu.br/periodicos/index.php/rpd/article/download/642/631/1254}
\item Record of the same topic: \url{https://oasisbr.ibict.br/vufind/Record/IFMT_f2dbaeecdcc990cadb914b9dd0885ec8}
\end{itemize}

\section{Linguistic and Reference Materials}
\begin{itemize}
\item \emph{Morfossintaxe da l\'ingua tapirap\'e} (doctoral thesis, UnB): \url{https://oasisbr.ibict.br/vufind/Record/UNB_a01804ab48b2793206ff05613764bc56}
\item Undergraduate thesis proposing a Tapirap\'e dictionary (UnB, 2019): \url{https://bdm.unb.br/bitstream/10483/27009/1/2019_NykMattheusCasseanoAndrade_tcc.pdf}
\item Tapirap\'e language (Wikipedia): \url{https://en.wikipedia.org/wiki/Tapirap%C3%A9_language}
\end{itemize}

\section{Religious and Community Resources}
None found in search results.

\section{Digital and Audio Resources}
No audio, video or software resources found in search results.

\section{Curriculum Development Context}
The language is still transmitted to children and used in community schools, with teachers coining neologisms for school vocabulary. The community regards Tapirap\'e as first language and Portuguese as second. A full dictionary and school textbooks were not found; the thesis and dictionary proposal are starting points, and the Tapi'it\~awa school and community should be consulted.
\end{document}